% open-logic.tex
% compiles to produce complete text using open-logic-dev style

\documentclass[../include/open-logic-part]{subfiles}

\begin{document}

\clearpage

\begin{editorial}
Berkas iki nglebokake kabeh isi kang kalebu ing Open Logic Project.
Cathetan panyunting kaya iki, yen katon, nuduhake yen berkas iki digawe
tanpa ngrancang carane materi iki bakal disuguhake. Yen bisa maca
cathetan iki, PDF iki bokmenawa \emph{ora prayoga} digunakake kanggo
mulang utawa sinau.

Open Logic Project nyedhiyakake akeh cara kanggo ngasilake teks kang
luwih trep kanggo piwulangan utawa sinau mandhiri. Tuladhane, kanthi
setelan baku, teks iki ndadekake kabeh operator logika minangka
operator primitif lan nggarap kabeh kasus kanggo saben operator ing
pambuktèn. Nanging, luwih becik sawetara kasus iki didadekake latihan.
Open Logic Project uga isih dikembangake. Kanggo nyengkuyung panyusunan
bebarengan lan pambecikan, materi tetep dilebokake sanajan mung awujud
draf, durung duwe latihan, lan sapanunggalane. PDF kang digawe kanggo
kuliah ora bakal ngemot perangan-perangan iki.

Kanggo nemokake PDF kang luwih trep kanggo mulang lan sinau, delengen
\href{http://builds.openlogicproject.org/}{PDF conto kanggo kuliah kang
  cumepak ing situs OLP}. Kanggo nggawe PDF dhewe, bisa miwiti saka
\href{https://github.com/OpenLogicProject/OpenLogic/tree/master/courses/sample}{conto
  berkas pangendhali} utawa ndeleng sumber buku ajar kang diasilake
kanggo conto kang luwih njlimet lan luwih maju.
\end{editorial}

\olimport[sets-functions-relations]{sets-functions-relations-complete}

\olimport[propositional-logic]{propositional-logic}

\olimport[first-order-logic]{first-order-logic}

\olimport[model-theory]{model-theory}

\olimport[computability]{computability}

\olimport[turing-machines]{turing-machines}

\olimport[incompleteness]{incompleteness}

\olimport[second-order-logic]{second-order-logic}

\olimport[lambda-calculus]{lambda-calculus}

\olimport[many-valued-logic]{many-valued-logic}

\olimport[normal-modal-logic]{normal-modal-logic}

\olimport[applied-modal-logic]{applied-modal-logic}

\olimport[intuitionistic-logic]{intuitionistic-logic}

\olimport[counterfactuals]{counterfactuals}

\olimport[set-theory]{set-theory}

\olimport[methods]{methods}

\olimport[history]{history}

\olimport[reference]{reference}

\end{document}
